% ============================================================
% CAPITOLO B - LE TRIADI
% ============================================================
\chapter{Le triadi}
\label{chap:triadi}

Dopo aver esaminato le fondamenta fisiche e gli intervalli nel Capitolo \ref{chap:fondamenti}, possiamo ora affrontare la costruzione degli accordi. L'elemento base dell'armonia tonale, la sua "molecola" fondamentale, è la \textbf{triade}. Quasi tutta la musica occidentale, dalla polifonia rinascimentale al pop moderno, si basa sulla successione e sulla manipolazione di questi accordi di tre note.

\section{Costruzione della triade}
\label{sec:costruzione_triade}

\begin{definizione}[La triade]
	\label{def:triade}
	Una \textbf{triade} è un accordo di tre note costruito sovrapponendo due intervalli di terza a partire da una nota base, detta \textbf{fondamentale}. Le altre due note prendono il nome di \textbf{terza} e \textbf{quinta} dell'accordo, in base all'intervallo che formano con la fondamentale.
\end{definizione}

La combinazione dei due tipi di terza (maggiore e minore) genera le diverse "specie" o "qualità" di triadi, ognuna con un colore e una funzione armonica distinti.

\begin{definizione}[Specie di triadi]
	\begin{itemize}
	\item \textbf{Triade maggiore:} terza maggiore + terza minore. È la triade "naturale", generata direttamente dalla serie armonica.
	\item \textbf{Triade minore:} terza minore + terza maggiore. Stabile, ma con un colore più scuro.
	\item \textbf{Triade diminuita:} terza minore + terza minore. Contiene il tritono (quinta diminuita), intervallo dissonante che genera forte tensione.
	\item \textbf{Triade eccedente (o aumentata):} terza maggiore + terza maggiore. L'intervallo di quinta aumentata conferisce un carattere ambiguo e instabile.
\end{itemize}
\end{definizione}

Di queste quattro specie, solo le prime due (maggiore e minore) sono considerate \textbf{consonanti} e possono fungere da punto di riposo. Le triadi diminuita ed eccedente sono \textbf{dissonanti} e richiedono risoluzione.

\begin{esempio}[Le quattro specie principali di triadi]
	\label{ex:quattro_triadi}
\musica{es_cap02_01}
\end{esempio}

% ============================================================
% INTEGRAZIONE SCIENTIFICA - FISICA
% ============================================================
\begin{fisica}[I rapporti di frequenza e la stabilità degli accordi]
	\label{phys:rapporti_frequenza_triadi}
	
	Per comprendere perché alcune triadi suonano più "stabili" di altre, dobbiamo analizzare i rapporti di frequenza tra le note che le compongono. Abbiamo visto nel Capitolo \ref{chap:fondamenti} che la serie armonica genera naturalmente rapporti semplici tra le frequenze.
	
	\textbf{La triade maggiore e i rapporti semplici}
	
	Consideriamo una triade di Do maggiore (Do-Mi-Sol). Se assumiamo che la fondamentale Do abbia frequenza $f_1$, le frequenze approssimative delle altre note nella scala temperata sono:
	\begin{itemize}
		\item Do (fondamentale): $f_1$
		\item Mi (terza maggiore): $\approx 1.26 \cdot f_1$ (rapporto ideale $5:4 = 1.25$)
		\item Sol (quinta giusta): $\approx 1.50 \cdot f_1$ (rapporto ideale $3:2 = 1.50$)
	\end{itemize}
	
	Nel sistema di \textbf{intonazione naturale} (o "giusta"), questi rapporti sono esatti: $4:5:6$. Questo significa che se Do vibra a 200 Hz, Mi vibra a 250 Hz e Sol a 300 Hz. Questi rapporti tra \textbf{numeri interi piccoli} sono ciò che il nostro orecchio percepisce come massimamente consonante.
	
	\textbf{La triade diminuita e i rapporti complessi}
	
	Confrontiamo ora con una triade diminuita (Do-\Mibemolle -\Solbemolle ). I rapporti approssimativi sono:
	\begin{itemize}
		\item Do: $f_1$
		\item \Mibemolle  (terza minore): $\approx 1.19 \cdot f_1$ (rapporto $6:5$)
		\item \Solbemolle  (quinta diminuita): $\approx 1.41 \cdot f_1$ (rapporto complesso, circa $\sqrt{2}:1$)
	\end{itemize}
	
	Il \textbf{tritono} (Do-\Solbemolle , rapporto $\sqrt{2}:1$ nel temperamento equabile) in particolare non può essere espresso come rapporto di numeri interi piccoli per la nota irrazionalità di $\sqrt{2}$. Se tentassimo di approssimarlo, otterremmo qualcosa come $45:32$ o $64:45$, numeri molto più grandi e complessi.
	
	\textbf{Il fenomeno dei battimenti acustici}
	
	Quando due frequenze vicine suonano insieme, si genera un fenomeno chiamato \textbf{battimento}. La frequenza del battimento è pari alla differenza tra le due frequenze:
	\[
	f_{\text{battimento}} = |f_2 - f_1|
	\]
	
	Con rapporti semplici (come 3:2), gli armonici delle due note tendono a coincidere o a essere molto distanziati, minimizzando i battimenti. Con rapporti complessi, invece, molti armonici cadono a frequenze vicine ma non coincidenti, generando battimenti multipli che l'orecchio percepisce come "ruvidità" o \textbf{dissonanza}.
	
	\textbf{Tavola comparativa}
	
	\begin{center}
	\begin{tabular}{|l|c|l|}
	\hline
	\textbf{Triade} & \textbf{Rapporti (ideali)} & \textbf{Percezione} \\
	\hline
	Maggiore & $4:5:6$ & Stabile, luminosa \\
	Minore & $10:12:15$ & Stabile, malinconica \\
	Diminuita & include tritono & Instabile, tesa \\
	Eccedente & include 5a aumentata & Ambigua, irrequieta \\
	\hline
	\end{tabular}
	\end{center}
	
	\textbf{Conclusione}
	
	La stabilità di un accordo ha quindi una base fisica oggettiva: triadi con rapporti di frequenza semplici generano meno interferenze acustiche e vengono percepite come più consonanti. Per questo le triadi maggiore e minore sono i pilastri dell'armonia tonale, mentre la triade diminuita richiede quasi sempre risoluzione verso accordi più stabili.
\end{fisica}


\section{Costruzione delle triadi sui gradi della scala}
\label{sec:triadi_sui_gradi}

Ogni grado della scala può essere la fondamentale di una triade. La qualità di questa triade è determinata dalla struttura della scala stessa. Analizziamo il modo maggiore.

\begin{definizione}[Triadi diatoniche nel modo maggiore]
	\label{def:triadi_diatoniche_maggiore}
	Le triadi costruite sui sette gradi della scala maggiore sono:
	\begin{itemize}
		\item \textbf{I, IV, V:} Triadi maggiori (funzione di stabilità)
		\item \textbf{II, III, VI:} Triadi minori (funzione secondaria)
		\item \textbf{VII:} Triade diminuita (funzione di tensione)
	\end{itemize}
\end{definizione}

\begin{esempio}[Le sette triadi diatoniche di Do Maggiore]
	\label{ex:triadi_diatoniche_do}
	\musica{es_cap02_02}
\end{esempio}

La nomenclatura dell'esempio \ref{ex:triadi_diatoniche_do} usa numeri romani maiuscoli per tutti gli accordi (I, II, III, IV, V, VI, VII), specificando tra parentesi la natura dell'accordo (M per maggiore, m per minore, dim per diminuito).

\section{I rivolti della triade}
\label{sec:rivolti_triade}

Finora abbiamo considerato le triadi nel loro \textbf{stato fondamentale}, ovvero con la fondamentale dell'accordo posizionata al basso. Tuttavia, è possibile disporre le note dell'accordo in modo diverso, ponendo al basso la terza o la quinta anziché la fondamentale. Queste diverse disposizioni sono chiamate \textbf{rivolti}.

\subsection{Basso fondamentale e basso reale}
\label{subsec:basso_fondamentale_reale}

Prima dei rivolti serve una distinzione terminologica.

\begin{definizione}[Basso fondamentale e basso reale]
	\label{def:basso_fondamentale_reale}
	In un accordo si distinguono due concetti:
	\begin{itemize}
		\item \textbf{Basso fondamentale:} la nota che dà il nome all'accordo e sulla quale l'accordo è costruito per terze sovrapposte. È un concetto teorico-funzionale, indipendente dalla disposizione effettiva delle note.
		\item \textbf{Basso reale:} la nota più grave effettivamente suonata, qualunque essa sia.
	\end{itemize}
	
	Se basso reale e fondamentale coincidono, si dice che l'accordo è in stato fondamentale, altrimenti si dice che è in stato di rivolto.
	
\end{definizione}

\begin{esempio}[Basso fondamentale vs. basso reale in Do Maggiore]
	\label{ex:basso_fond_reale}
	Lo stesso accordo di Do maggiore (basso fondamentale: Do) con tre diversi bassi reali.
\musica{es_cap02_03}
	In tutti e tre i casi l'accordo è sempre "Do maggiore" (basso fondamentale = Do), ma il basso reale cambia.
\end{esempio}


\subsection{Il primo rivolto delle triadi (stato di sesta)}
\label{subsec:primo_rivolto}

Il primo rivolto si ottiene ponendo la \textbf{terza} dell'accordo al basso.

\begin{definizione}[Primo rivolto delle triadi]
	\label{def:primo_rivolto}
	Un accordo è in \textbf{primo rivolto} quando la sua terza si trova al basso. Viene indicato con la cifratura \textbf{6} (o talvolta 6/3), poiché rispetto al basso la fondamentale si trova a distanza di sesta e la quinta a distanza di terza.
\end{definizione}

\begin{esempio}[Una triade in primo rivolto]
	\label{ex:primo_rivolto_triadi}
	\musica{es_cap02_04}
\end{esempio}

\subsection{Il secondo rivolto delle triadi (stato di quarta e sesta)}
\label{subsec:secondo_rivolto}

Il secondo rivolto si ottiene ponendo la \textbf{quinta} dell'accordo al basso.

\begin{definizione}[Secondo rivolto delle triadi]
	\label{def:secondo_rivolto}
	Un accordo è in \textbf{secondo rivolto} quando la sua quinta si trova al basso. Viene indicato con la cifratura \textbf{6/4}, poiché rispetto al basso la fondamentale si trova a distanza di quarta e la terza a distanza di sesta.
\end{definizione}

\begin{esempio}[Una triade in secondo rivolto]
	\label{ex:secondo_rivolto_triadi}
	\musica{es_cap02_05}
\end{esempio}

\subsection{La numerica: come si indica lo stato di un accordo}
\label{subsec:cifratura_accordi}

Lo stato di un accordo si scrive sotto il suo grado con una \textbf{numerica} (o \textbf{cifratura}). Per capire da dove nasce basta guardare l'accordo dal basso.

\begin{definizione}[Numerica di un accordo]
	\label{def:numerica_accordo}
	La \textbf{numerica completa} di un accordo è l'elenco degli intervalli che le altre note formano con il \textbf{basso reale}, contati a partire dal basso (senza badare alle ottave) e scritti in colonna, dal più grande in alto al più piccolo in basso. Non dipende dalla tonalità né dal grado della scala: dipende solo da quale nota dell'accordo sta al basso.
\end{definizione}

Poiché alcuni numeri sono sempre sottintesi, nella pratica si usa una \textbf{numerica sintetica}, che conserva le sole cifre necessarie a distinguere gli stati dell'accordo. La tabella confronta le due forme per le triadi e per le settime (queste ultime sono trattate dal Capitolo \ref{chap:settima_dominante}); gli esempi sono costruiti sull'accordo di Do maggiore e sull'accordo di settima di dominante Sol--Si--Re--Fa.

\begin{center}
\small\renewcommand{\arraystretch}{1.25}
\begin{tabular}{|l|l|c|l|c|}
\hline
\textbf{Stato} & \textbf{Basso} & \textbf{Completa} & \textbf{Sintesi possibili} & \textbf{Usata} \\
\hline
Triade, fondamentale & Do (fondamentale) & 5 3 & 5 3, 5, \textit{nessuna} & --- \\
Triade, primo rivolto & Mi (terza) & 6 3 & 6 3, 6 & 6 \\
Triade, secondo rivolto & Sol (quinta) & 6 4 & 6 4 & \cifre{6}{4} \\
\hline
Settima, fondamentale & Sol (fondamentale) & 7 5 3 & 7 5 3, 7 & 7 \\
Settima, primo rivolto & Si (terza) & 6 5 3 & 6 5 3, 6 5, 5 & \cifre{6}{5} \\
Settima, secondo rivolto & Re (quinta) & 6 4 3 & 6 4 3, 4 3 & \cifre{4}{3} \\
Settima, terzo rivolto & Fa (settima) & 6 4 2 & 6 4 2, 4 2, 2 & 2 \\
\hline
\end{tabular}
\end{center}

La numerica completa si legge direttamente sulle note: nel primo rivolto di Do maggiore (Mi--Sol--Do) sopra il Mi al basso ci sono il Sol, a distanza di terza, e il Do, a distanza di sesta; nel secondo rivolto (Sol--Do--Mi) sopra il Sol ci sono il Do (quarta) e il Mi (sesta). Nella settima di dominante con il Fa al basso (terzo rivolto) sopra di esso ci sono il Sol (seconda), il Si (quarta) e il Re (sesta): numerica completa $6\,4\,2$, sintetica 2.

\begin{regola}[Numerica usata nel manuale]
	\label{reg:numerica_manuale}
	Nel manuale ogni accordo è indicato dal suo \textbf{grado} (numero romano) con sotto la \textbf{numerica sintetica}:
	\begin{itemize}
		\item \textbf{triadi:} stato fondamentale, nessuna cifra; primo rivolto 6; secondo rivolto 6 e 4 (la 6 sopra la 4);
		\item \textbf{settime:} stato fondamentale 7 (la cifra serve a distinguere la settima dalla triade); primo rivolto 6 e 5; secondo rivolto 4 e 3; terzo rivolto 2.
	\end{itemize}
	Nelle figure il grado sta sopra e la numerica sotto, in colonna (per esempio il grado V con sotto 6 e 5, cifra sotto cifra). Nel testo corrente, per non alterare l'interlinea, la stessa sigla è scritta in forma compatta, con la prima cifra in alto e la seconda in basso: \sigla{I}{6}{}, \sigla{I}{6}{4}, \sigla{V}{7}{}, \sigla{V}{6}{5}, \sigla{V}{4}{3}, \sigla{V}{2}{}.
\end{regola}

La numerica sintetica non va confusa con la numerica del basso numerato, in cui altri numeri (ritardi come 4--3 o 9--8) e le alterazioni (\#, $\flat$, $\natural$) si aggiungono a questa: ne tratta il Capitolo \ref{chap:metodo_lavoro}.

\subsection{Il problema del secondo rivolto}
\label{subsec:problema_secondo_rivolto}

Il secondo rivolto differisce dallo stato fondamentale e dal primo rivolto: pur essendo costituito dalle stesse note di una triade consonante, viene percepito come \textbf{dissonante}. Il suo uso è quindi molto limitato.

\begin{fisica}[Il conflitto acustico dell'accordo di quarta e sesta]
	\label{phys:conflitto_6_4}
	
	Quando ascoltiamo un accordo, il nostro orecchio tende a identificare la nota più grave (il basso reale) come fondamentale "percepita". Questo avviene perché la nota al basso genera la propria serie armonica, che domina la percezione complessiva.
	
	\textbf{Il conflitto della quarta}
	
	Nel secondo rivolto, la fondamentale dell'accordo si trova a distanza di \textbf{quarta} dal basso. Questo intervallo è problematico per ragioni fisiche:
	
	\begin{itemize}
		\item La quarta (rapporto $4:3$) non compare tra i primi armonici semplici della serie del basso
		\item Si crea un'interferenza: l'orecchio "sente" il basso come fondamentale, ma la struttura dell'accordo suggerisce una nota diversa
		\item Questa ambiguità genera instabilità percettiva
	\end{itemize}
	
	\textbf{Rappresentazione matematica}
	
	Se il basso suona a frequenza $f_b$, i suoi armonici principali sono $f_b$, $2f_b$, $3f_b$, $4f_b$, $5f_b$... La quarta sopra il basso ha frequenza $f_4 = \frac{4}{3}f_b$. Questo rapporto non appare tra i primi armonici semplici, mentre la quinta ($3:2$) è il secondo rapporto più consonante dopo l'ottava.
	
	\textbf{Conclusione}
	
	L'accordo di 6/4 è dissonante per la disposizione delle note: l'orecchio percepisce il basso come fondamentale, mentre la struttura dell'accordo ne implica un'altra. Per questo l'accordo di quarta e sesta richiede sempre un trattamento particolare e non può essere usato liberamente come un rivolto qualsiasi.
\end{fisica}


